%Chargement des paquets
\usepackage{amsmath}
\usepackage{amsfonts}
\usepackage{amssymb}
\usepackage{enumerate}
\usepackage{mathtools}
\usepackage{bbm}
\usepackage{xparse, etoolbox}
\usepackage{enumerate}
\usepackage{enumitem}
\usepackage{mathabx}
\usepackage{babel}[french]

%environment
\usepackage{amsthm}
\usepackage{keytheorems}
\usepackage{hyperref} 

%Interface théorème
\newkeytheorem{lem}[
	name = Lemme
]

\newkeytheorem{prop}[
	name = Proposition
]

\newkeytheorem{thm}[
	name = Théorème
]

\newkeytheorem{coro}[
	name = Corollaire
]

\renewcommand*{\proofname}{Démonstration}

\theoremstyle{definition}
\newtheorem{defn}{Définition}

\theoremstyle{definition}
\newtheorem*{nota}{Notation}

\theoremstyle{definition}
\newtheorem*{limi}{Liminaire}

\theoremstyle{remark}
\newtheorem{rmq}{Remarque}

\theoremstyle{plain}
\newtheorem*{fait}{Fait}


%Ensembles de nombres
\newcommand{\N} {\mathbb{N}}
\newcommand{\Ne}{\N^\ast}
\newcommand{\Z} {\mathbb{Z}}
\newcommand{\Q} {\mathbb{Q}}
\newcommand{\R} {\mathbb{R}}
\newcommand{\Rb}{\overline{\mathbb{R}}}
\newcommand{\Rp}{\R_+}
\newcommand{\Rm}{\R_-}
\newcommand{\K} {\mathbb{K}}
\newcommand{\Ket} {\mathbb{K^{\ast}}}
\newcommand{\Cx}{\mathbb{C}}

%Ensembles, tribus
\newcommand{\A}{\mathbb{A}}
\renewcommand{\P}{\mathcal{P}}
\newcommand{\T}{\mathcal{T}}
\newcommand{\F}{\mathcal{F}}
\newcommand{\C} {\mathcal{C}}
\newcommand{\ouv}{\mathcal{O}}
\newcommand{\B}{\mathcal{B}}
\newcommand{\M}{\mathcal{M}}
\newcommand{\etag}{\mathcal{E}}
\newcommand{\etagOp}{\etag^{+}(\Omega)}
\newcommand{\etagO}{\etag(\Omega)}
%%Lp avant quotientage
\newcommand{\Lic}{\mathcal{L}^1}
\newcommand{\Lpc}{\mathcal{L}^p}
\newcommand{\Linfc}{\mathcal{L}^{\infty}}
\newcommand{\Lifull}{\Lic(\Omega, \B, m)}
\newcommand{\Lpfull}{\Lpc(\Omega, \B, m)}
\newcommand{\Linffull}{\Linfc(\Omega, \B, m)}
%%Lp après quotientage
\newcommand{\Li}{L^1}
\newcommand{\Ld}{L^2}
\newcommand{\Lp}{L^p}
\newcommand{\Linf}{L^{\infty}}
\newcommand{\Listd}{\Li(\Omega, m)} 
\newcommand{\Ldstd}{\Ld(\Omega, m)} 
\newcommand{\Lpstd}{\Lp(\Omega, m)} 

%Quelques opérateurs
\DeclareMathOperator{\codim}{codim}
\DeclareMathOperator{\rg}{rg}
\DeclareMathOperator{\tr}{tr}
\DeclareMathOperator{\vect}{Vect}
\DeclareMathOperator{\ima}{Im}
\DeclareMathOperator{\id}{Id}
\DeclareMathOperator{\sign}{sign}
\DeclareMathOperator{\ext}{ext}
\newcommand{\ind}{\mathbbm{1}}
\newcommand*{\spr}[2]{\left(\,#1\,\middle|\,#2\,\right)}
\newcommand{\interior}[1]{%
  {\kern0pt#1}^{\mathrm{o}}%
}
\DeclareMathOperator{\card}{card}
\DeclareMathOperator{\ppcm}{ppcm}

%Commandes de pure flemme
\newcommand{\veps}{\varepsilon}
\newcommand{\intab}{\int_{a}^{b}}
\newcommand{\intR}{\int_{-\infty}^{\infty}}
\newcommand{\intinf}{\int_{- \infty}^{+ \infty}}
\newcommand{\equi}{\Leftrightarrow}
\newcommand{\Kev}{$\K$-espace vectoriel }
\newcommand{\Kevs}{$\K$-espaces vectoriels }
\newcommand{\Rev}{$\R$-espace vectoriel }
\newcommand{\Revs}{$\R$-espaces vectoriels }
\newcommand{\Cev}{$\Cx$-espace vectoriel }
\newcommand{\Cevs}{$\Cx$-espaces vectoriels }
\newcommand{\evn}{(E, \norm{.})}
\newcommand{\interf}[2]{[#1,#2]}
\newcommand{\limb}{\lim\limits_}
\newcommand{\lebn}{\lambda_n}
\newcommand{\pp}{\text{~presque partout}}
\newcommand{\derivp}[2]{\frac{\partial #1}{\partial #2}}
\newcommand{\famiu}{(u_i)_{i \in I}}
\newcommand{\sev}{\text{~s.e.v.}}
\newcommand{\Mnp}{M_{n,p}(\K)}
\newcommand{\Mn}{M_{n}(\K)}
\newcommand{\tq}{\text{~t.q.~}}
\newcommand{\mq}{~montrer~que~}
\newcommand{\et}{\quad \text{et} \quad}
\newcommand{\ou}{\text{~ou~}}
\newcommand{\Kx}{\K[X]}


%Commandes de gars chelou
\renewcommand{\subset}{\subseteq}

%Commande restriction, ty duckduckgo
\def\restr#1#2{\mathchoice
              {\setbox1\hbox{${\displaystyle #1}_{\scriptstyle #2}$}
              \restraux{#1}{#2}}
              {\setbox1\hbox{${\textstyle #1}_{\scriptstyle #2}$}
              \restraux{#1}{#2}}
              {\setbox1\hbox{${\scriptstyle #1}_{\scriptscriptstyle #2}$}
              \restraux{#1}{#2}}
              {\setbox1\hbox{${\scriptscriptstyle #1}_{\scriptscriptstyle #2}$}
              \restraux{#1}{#2}}}
\def\restraux#1#2{{#1\,\smash{\vrule height .8\ht1 depth .85\dp1}}_{\,#2}} 

%Quelques normes
\DeclarePairedDelimiter\abs{\lvert}{\rvert}
\DeclarePairedDelimiter\labs{\bigl\lvert}{\bigr\rvert}
\DeclarePairedDelimiter\norm{\lVert}{\rVert}
\DeclarePairedDelimiterXPP\normi[1]{}\lVert\rVert{_1}{\ifblank{#1}{\:\cdot\:}{#1}}
\DeclarePairedDelimiterXPP\normd[1]{}\lVert\rVert{_2}{\ifblank{#1}{\:\cdot\:}{#1}}
\DeclarePairedDelimiterXPP\normp[1]{}\lVert\rVert{_p}{\ifblank{#1}{\:\cdot\:}{#1}}
\DeclarePairedDelimiterXPP\normq[1]{}\lVert\rVert{_q}{\ifblank{#1}{\:\cdot\:}{#1}}
\DeclarePairedDelimiterXPP\norminf[1]{}\lVert\rVert{_\infty}{\ifblank{#1}{\:\cdot\:}{#1}}

%Bracket en angle
\DeclarePairedDelimiter\angl{\langle}{\rangle}

%Set, parenthèses
\newcommand{\set}[1]{\{ #1 \}}
\newcommand{\bigset}[1]{\bigl\{ #1 \bigr\}}
\newcommand{\Bigset}[1]{\Bigl\{ #1 \Bigr\}}
\newcommand{\bigpar}[1]{\bigl( #1 \bigr)}
\newcommand{\Bigpar}[1]{\Bigl( #1 \Bigr)}

%Opitmisation des opérateurs de comparaison
\DeclarePairedDelimiter\tio{o (}{)}
\DeclarePairedDelimiter\gro{O (}{)}


\pagestyle{fancy}

\fancyhead[C]{Théorème de Stone-Weierstrass}
\fancyhead[R]{}

\begin{document}

\underline{Source :} Rombaldi - Analyse - pages 247-248

\underline{Leçons :}
\begin{itemize}
\item 201 : Espaces de fonctions. Exemples et applications.
\item 203 : Utilisation de la notion de compacité.
\item 209 : Approximation d’une fonction par des fonctions régulières. Exemples d’applications.
\end{itemize}

Dans ce développement, $(E,d)$ est un espace métrique que l'on désignera simplement par $E$.	

\begin{defn}[Complétude]
On dit qu'une partie $K$ de $E$ est complète si toute suite de Cauchy d'éléments de $K$ converge dans $K$.
\end{defn}

\begin{defn}[Précompacité]
On dit qu'une partie $K$ de $E$ est précompacte si, pour tout réel $r>0$, 
$K$ peut être recouvert par un nombre fini de boules ouvertes de rayon $r$ et centrées en des éléments de $K$. 
\end{defn}

\begin{defn}[Compacité]
On dit qu'une partie $K$ de $E$ est compacte si, de toute suite de points de $K$, on peut extraire une sous-suite qui converge 
dans $K$.
\end{defn}

\begin{lem}
\label{lem:precomp}
Toute partie compacte de $E$ est précompacte.
\end{lem}

\begin{proof}
Supposons qu'il existe $r>0$ tel que $K$ ne puisse pas être recouvert par un nombre \textbf{fini} de boules ouvertes de rayon $r$.
Prenons $u_0 \in K$, par hypothèse, il existe $u_1 \in K \setminus B(u_0,r)$.
De proche en proche, on construit une suite $(u_n)_{\N}$ telle que
\[ \forall n \geq 1, u_{n+1} \in K \setminus \bigcup_{k=0}^{n} B(u_k, r). \]
Par compacité, on peut extraite le suite $(u_n)$ une suite $(u_n)_S$ convergente ($S$ est une partie infinie de $\N$). 
En particulier, cette sous-suite est de Cauchy, donc il existe un rang $N \in S$ tel que :
\[ \forall n \in S : n, m \geq N \Rightarrow d(u_n,u_m) < r \]
ce qui est contradictoire, par construction.
\end{proof}

\begin{lem}
Toute partie compacte de $E$ est fermée.
\end{lem}

\begin{proof}
Soit $(u_n)_S$ une suite convergente de $K$, comme il existe une suite extraite de $(u_n)$ convergente dans $K$ 
et qu'une suite extraite a même limite que sa suite parente, on en déduit que $(u_n)_S$ converge dans $K$.
\end{proof}

\begin{prop}
Toute partie compacte de $E$ est séparable.
\end{prop}

\begin{proof}
Soit $K$ une partie compacte de $E$, selon le lemme $\ref{lem:precomp}$, 
il existe un recouvrement fini de $K$ par des boules de rayon $r>0$ arbitraire.
On note 
\[ \B=\{ B(u(n,k), 1/2^n) ; n \in \N, 1 \geq k \geq m_n \}, \]
avec, pour $n \in \N$ fixé, les boules $B(u(n,k), 1/2^n)$ recouvrant l'ensemble $K$ ($m_n$ étant le nombre de ces boules pour le rayon $1/2^n$).
$\B$ est dénombrable, comme réunion dénombrable d'ensembles finis, par suite 
\[ U = \{ u(n,k) ; n \in \N, 1 \leq k \leq m_n \}  \subset K \]
est un ensemble dénombrable d'éléments de $K$, il reste à montrer qu'il est partout dense dans $K$. Soit $\veps >0$ et $x \in K$, on peut toujours prendre $n \in \N$, tel que $1/2^n < \veps$, puis, par construction des recouvrements, on prend $k, 1 \leq k \leq m_n$ tel que 
\[ x \in B(u(n,k), 1/2^n) . \]
Ainsi, $d(x, u(n,k)) < 1/2^n < \veps$, et $K$ est bien séparable.
\end{proof}

\begin{lem}[Lebesgue]
\label{lem:lebesg}
Soient $K$ un compact de $E$ et $(O_i)_{i \in I}$ une famille d'ouverts de $E$ recouvrant $K$. Alors, il existe un réel $r>0$ tel que toute boule ouverte de centre $x \in K$ et de rayon $r$ soit contenue dans l'un des $O_i$.
\end{lem}

\begin{proof}
Supposons que cette propriété soit fausse, alors :
\[ (\forall r > 0) : (\exists x_r \in K ; \forall i \in I, B(x_r, r) \not \subset O_i). \]
En particulier, en prenant pour $r$ les valeurs $1/2^n$ pour $n \in \N$, on construit une suite $(x_n)_{\N}$ d'éléments de $K$ vérifiant 
\[ \forall n \in \N, \forall  i \in I, B(x_n, 1/2^n) \not \subset O_i). \]
Par compacité de $K$, on peut extraire de cette suite une sous-suite de support $S \subset \N$ (S partie infinie) et convergente. Notons $x$ la limite de cette suite, par hypothèse, il existe $i \in I$ tel que $x \in O_i$ et, par suite, il existe $r >0$ tel que $B(x,r) \subset O_i$. Prenons $\veps \in ]0, r/2 [$, il existe un rang $N \in S$ tel que 
\[ \forall n \in S : n \geq N \Rightarrow d(x_n,x) < \veps . \]
Considérons maintenant $x' \in B(x_n, r/2)$, on a
\[ d(x', x) \leq d(x', x_n) + d(x_n, x) < r \]
donc $\forall n \in S : n \geq N \Rightarrow B(x_n, r/2) \subset O_i$. En prenant $n \in S, n \geq N$ tel que $1/2^n < r/2$ on arrive à une contradiction.
\end{proof}

\begin{lem}
\label{lem:adhCauchy}
Soit $(u_n)_S$ une suite de Cauchy de $E$. Si $(u_n)_S$ admet une valeur d'adhérence, alors $(u_n)_S$ est convergente.
\end{lem}

\begin{proof}
Notons $u$ la valeur d'adhérence de cette suite et considérons $\veps >0$. Prenons $N \in S$ tel que 
\[ \forall p, q \in S, p \geq N, q \geq N, d(u_p, u_q) < \veps/2 , \]
Comme il existe $n_0 \in S, n_0 \geq N$ tel que $d(u_{n_0}, u) < \veps/2$ on conclut que 
\[ \forall n \in S : n \geq n_0 \Rightarrow d(u_n, u) < \veps \]
ce qui démontre la convergence de $(u_n)_S$ vers $u$.
\end{proof}

\begin{thm}
Soit $K$ une partie non vide de $E$. Les propriétés suivantes sont équivalentes :
\begin{enumerate}[label=(\roman*)]
\item $K$ est compact ;
\item de tout recouvrement de $K$ par des ouverts, on peut en extraire un sous-recouvrement fini (propriété de Borel Lebesgue) ;
\item $K$ est précompact et complet.
\end{enumerate}
\end{thm}

\begin{proof}
$(i) \Rightarrow (ii)$ Supposons que $K$ est compact et qu'il existe une famille $(O_i)_{i \in I}$ d'ouverts de $E$ telle que 
\[ K \subset \bigcup_{i \in I} O_i . \]
Selon le lemme $\ref{lem:lebesg}$, il existe un rayon $r>0$ tel que, 
toute boule ouverte de centre $x \in K$ et de rayon $r$ soit contenue dans l'un des $O_i$. Or
\[ \bigcup_{x \in K} B(x, r) \]
forme un recouvrement trivial de $K$, par précompacité de $K$ on extrait de ce recouvrement un recouvrement fini noté 
\[ \bigcup_{k=1}^n B(x_k, r), n \in \N . \]
Comme chacune des boules du recouvrement appartient à un ouvert de la famille, on conclut aisément. \newline

$(ii) \Rightarrow (iii)$ Pour $r>0$, on considère le recouvrement :
\[ K \subset \bigcup_{x \in K} B(x,r) \]
dont on extrait un sous-recouvrement fini, ce qui démontre la précompacité de $K$. 
Supposons maintenant qu'il existe une suite $(u_n)_S$ de $K$ n'admettant pas de valeur d'adhérence. Alors
\[ \forall x \in K, \exists \veps_x >0, \exists N_x \in S, \forall n \geq N_x, d(u_n, x) \geq \veps_x . \]
Prenons le recouvrement 
\[ K \subset \bigcup_{x \in K} B(x, \veps_x) \]
dont on extrait un sous recouvrement fini. $n \geq 1$ étant le cardinal de ce recouvrement, on note
 $x_1, \dots x_n \in K$ et $\veps_1, \dots,\veps_n \in ]0, \infty[$ les centres et rayons tels que 
\[ K \subset \bigcup_{i =1}^n B(x_i, \veps_i) . \]
Notons $N = \max \{ N_x, x \in \{x_1, \dots, x_n \} \}$, alors, pour $n \in S, n \geq N$, 
le terme $u_n$ n'appartient pas à $\bigcup_{i =1}^n B(x_i, \veps_i)$, ce qui est contradictoire. 
Donc, toute suite de $K$ admet une valeur d'adhérence. 
On conclut en utilisant le lemme $\ref{lem:adhCauchy}$. \newline

$(iii) \Rightarrow (i)$ Soit $(u_n)_S$ une suite de poins de $K$. Notons, pour $N \in \N$
\[ A_N = \{u_n ; n \in S, n \geq N \} \text{ et } \mu_N = \inf \{ d(u_N, u_p) ; p \in S, p \geq N \} ,\]
Supposons que $\mu_N >0$, on recouvre $K$ par un nombre fini de boules de rayon $\mu_N$, 
alors au moins deux points de l'ensemble $A_N$ sont dans la même boule, ce qui est contradictoire. 
Donc $\mu_N=0$. Autrement dit
\[ \forall \veps > 0, \forall N \in \N, \exists p \geq N, \tq d(u_p, u_N) < \veps. \]
On construit donc une suite extraite $(u_n)_T$ ($T \subset S$) telle que :
\[ \forall N \in T, p \geq N, q \geq N  \Rightarrow d(u_p, u_q) < 1/2^{n+1} \]
autrement dit une suite de Cauchy. Cette suite converge dans $K$ par complétude, $K$ est donc compact.

\end{proof}

\begin{lem}
\label{lem:Weier}
La suite de fonctions polynomiales $(P_n)_{n \in \N}$, définie sur $[-1,1]$ par $P_0(x)=0$ et $P_{n+1}(x)=P_n(x) + \frac{x^2-P_n^2(x)}2$ converge uniformément sur $[-1,1]$ vers la fonction $x \mapsto \abs{x}$.
\end{lem}

\begin{proof}
On commence par \mq 
\[ \forall n \in \N : 0 \leq P_n(x) \leq P_{n+1}(x) \leq \abs{x} \]
Comme $P_1(x)=x^2/2$, l'affirmation est vraie pour $n=0$. On va la supposer vraie pour $n \in \N$ quelconque. On a 
\[ P_{n+2}(x) = P_{n+1}(x) + \dfrac{x^2 - P_{n+1}^2(x)}2 . \]
Or, par croissance de la fonction $x \mapsto x^2$ sur $\Rp$, on a $0 \leq P_{n+1}^2(x) \leq x^2$ et donc 
\[ P_{n+2}(x) - P_{n+1}(x) \geq 0 \qquad \text{et} \qquad P_{n+2}(x) \geq 0 . \]
De plus, $\abs{x} - P_{n+2}(x) = (\abs{x} - \frac12 x^2) - ( P_{n+1}(x) - \frac12 P_{n+1}^2(x)$. 
Par croissance de la fonction $x \mapsto x - \frac12 x^2$ sur [0,1] on conclut que 
\[ \abs{x} - P_{n+2}(x) \geq 0 \text{~~car~} 0 \leq P_{n+1}(x) \leq \abs{x} \leq 1 . \]
L'affirmation est donc vraie pour $n+1$, elle est donc vraie pour tout $n \in \N$. \newline

On a donc montré que, pour tout $x \in [-1, 1]$, la suite $(P_n(x))_{n \in \N}$ était croissante et majorée, elle admet donc une limite l(x).
Cette limite vérifie $l(x)=l(x)+\frac{x^2-l(x)^2}2$ soit $x^2-l(x)^2=0$ et donc $l(x)=\abs{x}$. 
La suite de fonctions $(P_n)_{n \in \N}$ est croissante dans $\C^0([-1,1], \R)$ 
et converge simplement vers la fonction continue $x \mapsto \abs{x}$ donc,
d'après le théorème de Dini, elle converge uniforémement vers cette fonction.
\end{proof}

\begin{defn}[Algèbre]
On appelle $\K$-algèbre un \Kev $\A$ muni d'une seconde loi interne (notée $\times$) vérifiant :
\begin{enumerate}[label=(\roman*)]
\item $(\A, +, \times)$ est un anneau ;
\item $\forall (x,y) \in \A^2, \forall \lambda \in \K, \lambda \cdot (x \times y) = (\lambda \cdot x) \times y = x \times (\lambda \cdot y)$.
\end{enumerate}
Si de plus, $(\A, +, \times)$ est munie d'un élément neutre pour la loi $\times$, on parle d'algèbre unitaire.
\end{defn}

\begin{defn}[Algèbre normée]
Soit $\A$ une $\K$-algèbre uniatire d'unité notée 1. Une norme $N$ sur $\A$ est dite sous-multiplicative si elle vérifie :
\[ \forall (x,y) \in \A^2, N(xy) \leq N(x)N(y). \]
Une algèbre normée est une algèbre munie d'une norme sous-multiplicative $N$ et telle $N(1)=1$. 
Dans le cas où de plus l'espace normé $(\A, N)$ est complet, on parle d'algèbre de Banach.
\end{defn}

\begin{thm}[Stone-Weierstrass]
Soit $K$ un compact contenant au moins deux points dans un espace métrique $E$ et 
$\A$ une sous-algèbre de $\C^0(K, \R)$ qui contient les constantes et sépare les points de $K$
dans le sens où pour tous $x \neq y$ dans $K$, il existe $f \in \A$ \tq $f(x) \neq f(y)$.
Dans ce cas $\A$ est partout dense dans $(\C^0(K, \R), \norminf{.})$.
\end{thm}

\begin{proof}
On divise la démonstration en plusieurs étapes pour la rendre plus digeste.
\begin{enumerate}[start=1,label={\bfseries Etape~\arabic* :}, wide, labelwidth=!, labelindent=0pt]
\item
On va d'abord \mq, pour toute fonction $f \in \A$, la fonction $\abs{f}$ est dans l'adhérence de $\A$ (relativement à $\C^0(K, \R)$). 
C'est clair pour la fonction nulle. 
Si $f \neq 0$, la fonction $x \mapsto \frac{f}{\norminf{f}}$ est bien définie et à valeurs dans $[-1,1]$.
On considère la suite de fonctions polynomiales $(P_n)_{n \in \N}$ définie dans le lemme $\ref{lem:Weier}$ et,
pour $n \in \N$, on définit la fonction $g_n = P_n \left ( \frac{f}{\norminf{f}} \right )$. 
La suite $(g_n)_{n \in \N}$ est une suite de $\A$ puisque $\A$ est une sous-algèbre de $\C^0(K, \R)$ qui contient les constantes.
De plus, cette suite converge uniformément vers la fonction $g = \frac{\abs{f}}{\norminf{f}}$ car
\[ \abs{g_n(x) - g(x)} = \abs*{ P_n \left ( \frac{f}{\norminf{f}} \right ) - \frac{\abs{f}}{\norminf{f}}} \leq 
\sup_{t \in [-1;1]} (P_n(x) - \abs{x}) , \]
car les fonctions $g_n, n \in \N$ et $g$ sont à valeurs dans $[-1,1]$ et
\[ \forall \veps >0, \exists N \in \N ; \forall n \in \N : n \geq N \Rightarrow \sup_{t \in [-1;1]} (P_n(x) - \abs{x}) < \veps , \]
toujours selon le lemme $\ref{lem:Weier}$. En d'autres termes, la suite de fonctions $g_n$ converge vers $g$ pour la norme infinie, 
soit $g \in \bar A$ et, par suite, $f = \norminf{f} g \in \bar A$.
\item
Il en résulte que, pour toutes fonctions $f$ et $g$ dans $\A$, 
les fonctions $\max(f,g)=\frac{f+g}2 + \frac{\abs{g-f}}{2}$ et $\min(f,g)=\frac{f+g}2 - \frac{\abs{g-f}}{2}$ appartiennent à $\bar \A$.
On en déduit (par récurrence) que pour toute suite $(f_k)_{1 \leq k \leq n}$ de fonctions de $\A$, 
les fonctions $\max_{1 \leq k \leq n} f_k$ et $\min_{1 \leq k \leq n} f_k$ sont dans $\bar \A$.
\item
On vérifie ensuite que pour deux points $x \neq y$ dans $K$ et pour deux réels quelconques $a$ et $b$,
il existe une fonction $g \in \A$ telle que $g(x)=a$ et $g(y)=b$. 
Par hypothèse, il existe $f \in \A \tq f(x) \neq f(y)$, on définit donc $g$ sur $K$ par :
\[ g(z) = a\dfrac{f(z) - f(y)}{f(x) - f(y)} + b\dfrac{f(z) - f(x)}{f(y) - f(x)} . \]
\item 
On prend une fonction $f \in \C^0(K, \R)$ et $\veps >0$. 
On va \mq il existe une fonction $g \in \bar A$ telle que $\norminf{f-g} < \veps$.
Pour $x \in K$ donné et pour tout $y \neq x$ dans $K$, selon l'étape 3, il existe une fonction $g_y \in \A$ telle que
$g_y(x)=f(x)$ et $g_y(y)=f(y)$. Pour chacun des $y$ considérés, on définit l'ensemble 
\[ O_y = \set{t \in K ; g_y(t) > f(t) - \veps } . \]
Ces ensembles sont des ouverts de $E$ (comme préimages par une fonction continue d'ouverts de $\R$) 
et ils contiennent les points $x$ et $y$. De plus 
\[ K \subset \bigcup_{y \in K \setminus \set{x}} O_y, \]
par compacité de $K$, on peut donc extraire un sous-recouvrement fini de ce recouvrement que l'on note 
\[ K \subset \bigcup_{1 \leq i \leq N} O_{y_i} , \]
les points $y_i$ étant tous différents de $x$.
Selon l'étape 2, la fonction $f_x = \max_{1 \leq i \leq N} g_{y_i}$ est dans $\bar A$ et elle vérifie $f_x(x)=f(x)$. De plus 
\[ \forall t \in K, \exists i \in \set{1, \dots, N} ; t \in O_{y_i} , \]
et on a 
\[ f_x(t) \geq g_{y_i}(t) > f(t) - \veps . \]
Pour tout $x \in K$, l'ensemble 
\[ \Omega_x = \set{t \in K ; f_x(t) < f(t) + \veps } \] 
est un ouvert de $E$ (toujours par continuité) qui contient x. Du recouvrement
\[ K \subset \bigcup_{x \in K} \Omega_x \]
on extrait un sous recouvrement fini 
\[ K \subset \bigcup_{1 \leq i \leq M} \Omega_{x_i} . \]
On définit ensuite la fonction $g = \min_{1 \leq i \leq M} f_{x_i}$ qui est dans $\bar A$ selon l'étape 2.
Pour tout $t \in K$, il existe un entier $i$ compris entre 1 et $M$ et tel que $t \in \Omega_{x_i}$ et on a :
\[ g(t) \leq f_{x_i}(t) < f(t) + \veps \]
Comme de plus, $\forall t \in K, \forall i \in \set{1, \dots, M}, f(t) - \veps < f_{x_i}$ on conclut que :
\[ \forall t \in K :  f(t) - \veps < g(t) <  f(t) + \veps . \]
Donc $\norminf{g-f} < \veps$, ce qui démontre que l'algèbre $\bar A$ est partout dense dans $(\C^0(K, \R), \norminf{.})$.
Il en est donc de même pour l'algèbre $A$.
\end{enumerate}
\end{proof}

\begin{coro}
Toute fonction continue $f:[a,b] \to \R$ est limite uniforme sur $[a,b]$ d'une suite de fonctions polynomiales.
\end{coro}

\begin{proof}
$[a,b]$ est un compact de $(\R, \abs{.}$ contenant au moins deux points.
L'ensemble des fonctions polynomiales sur $[a,b]$ est bien une algèbre séparante.
On applique directement le théorème précédent.
\end{proof}

\end{document}